\paragraph{A vertex-closure construction.}
Consider the two-regime model with $M=1$, $p_L=0$, $p_H=1$, zero background
switching, zero duration below the root, and taxon-specific deterministic
inheritance kernels. For $z\in\{0011,0101,0110\}$, let $j(z)$ be the
corresponding quartet resolution $12|34$, $13|24$, or $14|23$.

\begin{proposition}
For a root branch of finite duration $T\geq0$,
\[
 P(Q_i\mid z)=\frac{e^{-2T}}3+(1-e^{-2T})\mathbf 1\{i=j(z)\}.
\]
Consequently, every simplex vertex lies in the closure of the image of the
unrestricted two-regime model.
\end{proposition}

\begin{proof}
Under the specified pattern, exactly two background-homogeneous pairs are
available and each pair coalesces at rate one. Their independent survival
probability through the root branch is $e^{-2T}$. If a pair coalesces, its
quartet resolution is $j(z)$. If neither pair coalesces, the unresolved
ancestral quartet has probability $1/3$ for each topology. Combining these
events gives the formula. Letting $T\to\infty$ gives the vertex.
\end{proof}

Under finite branch durations and finite generator rates, a positive
four-lineage unresolved mass entering the root has positive probability for
each absorbing topology, because the residual vectors assign positive mass to
every unresolved background state. Hence such a finite-parameter calculation
cannot attain a simplex vertex. The construction reaches the vertices only in
the $T\to\infty$ closure.

This proposition does not imply full-simplex coverage, density of the full
simplex, or validity under identity or founding inheritance. The original
finite-$M$ Moran model is a different count-state process and is not covered
by the proposition.
